\subsection{应用于量子化}\label{scbq}

现在讨论如何把全纯极化用于量子化。考虑第 \ref{amo} 节引入的情形：$Y$ 是复辛流形，复结构为 $I$，全纯辛形式为 $\Omega$，并带反全纯对合 $\tau$ 使 $\tau^*(\Omega)=\bar\Omega$；$M$ 是 $\tau$ 不动点集的一个分支，对 $\omega_J=\Re\Omega$ 为辛的（同时对 $\omega_K=\Im\Omega$ 必然为拉格朗日的）。假定 $Y$ 有全纯极化 $\varPi$，其叶是某全纯映射 $\pi:Y\to W$ 的纤维，$W$ 为某复流形。我们要理解 $\varPi$ 的存在对 $M$ 的膜量子化意味着什么。

由于 $Y$ 被 $\varPi$ 的叶分层，每点 $p\in M$ 含于唯一的叶 $F_p$。$F_p$ 对 $\Omega$ 为拉格朗日，故对 $\omega_J$ 也是。于是 $\omega_J$ 限制到 $F_p\cap M$ 为零。因此 $F_p\cap M$ 在 $M$ 中至多是中间维的。若它是中间维的，它就是 $M$ 的拉格朗日子流形。假定对每个与 $M$ 相交的 $\varPi$ 叶都如此，且所有这些拉格朗日子流形都是 $\R^n$ 的拷贝。那么形如 $F_p\cap M$ 的 $M$ 的拉格朗日子流形定义 $M$ 的一个实极化 $\P$。在此情形，我们将证明 $M$ 的膜量子化等价于用实极化 $\P$ 作几何量子化。

设 $N$ 是 $\P$ 的叶空间。它是 $W$ 中参数化与 $M$ 相交的 $\varPi$ 叶的子空间。$N$ 中的点对应 $\varPi$ 的叶 $F$，使 $F\cap M$ 的实维数为 $n$，故 $F\cap M$ 在 $F$ 中是中间维的。$F\cap M$ 是全实的（因 $M$ 是全实的），故 $F$ 是 $F\cap M$ 的复化。由于 $F\cap M$ 为 $\tau$-不变，而 $F$ 是 $F\cap M$ 的复化，$F$ 被 $\tau$ 映到自身（虽非逐点）。

再由维数计数可知 $N$ 在 $W$ 中是中间维的；$N$ 的实维数为 $n=\frac12\dim_\R M=\frac12\dim_\C Y$，与 $W$ 的复维数相同。$N$ 参数化 $W$ 的一个中间维全实子空间，故 $W$ 可看作 $N$ 的复化。$\tau$ 作为 $W$ 的反全纯对称作用，保持 $N$ 不动。$W$ 中不属于 $N$ 的点对应不与 $M$ 相交的 $\varPi$ 叶；至少在 $N$ 的邻域内，这样的叶被 $\tau$ 成对交换。特别地，全纯极化 $\varPi$ 是 $\tau$-不变的，尽管并非每片叶都被 $\tau$ 映到自身。

一个初等例子：取 $Y=\C^2$，用复变量 $p,q$ 参数化，$\Omega=\d p\,\d q$。令 $\tau$ 由 $(p,q)\to(\bar p,\bar q)$ 作用，不动点集为 $M$。把 $p,q$ 按实部与虚部展开：$p=p_1+\i p_2$，$q=q_1+\i q_2$。$M$ 是由 $p_1,q_1$ 参数化的 $\R^2$ 拷贝。$Y$ 有全纯极化 $\varPi$，其叶为 $q=b$，$b$ 为复参数，标志叶的选择。于是叶空间 $W$ 是 $\C$ 的拷贝。$\tau$ 把 $q=b$ 的叶映到由 $\bar q=b$ 即 $q=\bar b$ 定义的叶；$\tau$-不变的叶是 $b$ 为实的叶。故 $\tau$ 通过复共轭作用在 $W\cong\C$ 上，固定实轴 $N$。$N$ 中的点对应与 $M$ 交于 $\R$ 的拷贝的 $\varPi$ 叶；这些 $\R$ 的拷贝就是 $M$ 的实极化 $\P$ 的叶。

一个稍精致的例子：考虑式 (\ref{nuuxx}) 的复辛流形 $Y'$，取 $j^2<0$，使 $s$ 为实。$M$ 由同一方程且 $x,y,z$ 为实描述。式 (\ref{horiz}) 给出 $Y'$ 的两个不同全纯极化 $\varPi$ 与 $\varPi'$，一个由 $q$ 参数化，一个由 $q'$ 参数化。$\tau$ 把 $\varPi$ 映到自身，也把 $\varPi'$ 映到自身。与 $M$ 有非空交集的 $\varPi$ 或 $\varPi'$ 的叶，是 $q$ 或 $q'$ 为实的叶。这些叶给出 $M$ 的两个实极化 $\P$ 与 $\P'$。与 $\varPi$、$\varPi'$ 相应的两个投影 $\pi,\pi':Y\to W\cong\bCP^1$ 限制到 $M$ 上，给出投影 $\pi,\pi':M\to N=\RP^1$，$N$ 由 $q$ 或 $q'$ 加上一个无穷远点参数化。

现在考虑相当不同的可能性。若 $F$ 是全纯极化 $\varPi$ 的叶，则 $F$ 与 $M$ 在 $Y$ 中都是中间维的，故 $F\cap M$ 自然可能只由一点组成（由若干点组成也自然，但我们不考虑这种情形）。假定对每个 $\varPi$ 叶，$F\cap M$ 至多是一点。于是投影 $\pi:Y\to W$ 限制为 $M$ 到其像 $\pi(M)$ 的同构（$\pi(M)$ 可能只是 $W$ 的一部分，后面例子会看到）。由于 $W$ 是复流形，这给 $M$ 一个复结构，记作\footnote{此记号的理由是：若 $Y$ 的复结构（我们称为 $I$）可延拓为带附加复结构 $J,K$ 的超 Kähler 结构，则 $\J$ 有时就是 $J$ 在 $M$ 上的限制。}$\J$。在复结构 $\J$ 下，$M$ 上的全纯函数是 $W$ 上全纯函数经 $\pi^*$ 的拉回。注意：若 $V$ 是 $Y$ 上切于 $\varPi$ 叶的向量场（不必全纯），则 $V$ 在 $W$ 上投影为零，故从 $W$ 拉回的函数的 $V$ 方向导数为零：
\be\label{looj} \mathcal L_V\cdot\pi^*(f)=0. \ee
在此情形，$\J$ 是几何量子化意义下 $M$ 的复极化，因为 $\omega_J|_M=\Re\Omega|_M$ 对 $\J$ 为 $(1,1)$ 型。为看出这点，令 $T_pM$ 为 $M$ 在 $p\in M$ 处的切空间，$T_{p,\C}M=T_{p,M}\otimes_\R\C$ 为其复化。考虑分解 $T_{p,\C}M=T^{1,0}_{p,\C}M\oplus T^{0,1}_{p,\C}M$，两个子空间分别是 $M$ 的全纯与反全纯复化切方向（在复结构 $\J$ 下）。作为向量空间，$T^{0,1}_{p,\C}M$ 可自然等同于 $T_pF_p$，即 $M$ 在含 $p$ 的 $\varPi$ 叶 $F_p$ 中的切空间（这由式 (\ref{looj}) 推出：$M$ 上的全纯函数被 $T_pF_p$ 方向的导数零化）。由于 $F_p$ 对 $\Omega$ 为拉格朗日，这说明 $\Omega|_M$ 的 $(0,2)$ 部分在复结构 $\J$ 下为零。连同 $M$ 对 $\Im\Omega$ 为拉格朗日，可知 $\Re\Omega|_M$ 的 $(2,0)$ 与 $(0,2)$ 部分都为零，故 $\Re\Omega|_M$ 为 $(1,1)$ 型。

$Y$ 的反全纯对合 $\tau$ 把 $\varPi$ 映到全纯极化 $\varPi'=\tau(\varPi)$。$\tau$ 把全纯变化的 $\varPi$ 叶族映到反全纯变化的 $\varPi'$ 叶族，故它反转 $M$ 上诱导的复结构。特别地，$\varPi$ 与 $\varPi'$ 总是不同。若 $\omega$ 在 $\varPi$ 诱导的复结构 $\J$ 下正定（即 $\J$ 给出几何量子化意义下的复 Kähler 极化），则它在 $\varPi'$ 诱导的复结构 $-\J$ 下负定。

在此情形，我们将在第 \ref{branecp} 节论证：$M$ 的膜量子化（至少作为向量空间，暂且不管 Hilbert 空间结构）与用 $\varPi$ 诱导的 $M$ 的复极化作几何量子化一致。Hilbert 空间结构一般不同，第 \ref{flatg} 节的例子会说明这一点。

但先看几个由 $Y$ 的全纯极化诱导 $M$ 的复极化的例子。初等例子仍取 $Y=\C^2$，$p,q,\Omega,\tau,M$ 同前。但现在用方程 $q+\i p=b$ 定义 $\varPi$ 的叶，$b$ 仍为标志叶选择的常数。这样的叶与 $M$ 交于唯一点 $q_1+\i p_1=b$。故此时 $M$ 获得一个复结构，其中 $z=q_1+\i p_1$ 全纯。切于叶 $q+\i p=b$ 的向量场是 $\frac{\partial}{\partial q}+\i\frac{\partial}{\partial p}$，沿 $M$ 解释为 $2\frac{\partial}{\partial\bar z}$。$\varPi$ 的 $\tau$-共轭是极化 $\varPi'$，其叶由 $q-\i p=b$ 定义，$b$ 仍为参数。诱导的复结构是使 $\bar z=q_1-\i p_1$ 全纯的那个。

稍微妙的例子：考虑第 \ref{example} 节研究过的复流形 $x^2+y^2+z^2=\j^2$，取 $\j^2>0$。$M$ 是 $x,y,z$ 为实的子流形。对实 $x,y,z$，式 (\ref{pok}) 的矩阵 $T$ 是 Hermite 的。特征值为 $\j,-\j$ 的 $2\times2$ Hermite 矩阵在给定一个特征向量时唯一确定，而可通过给定 $q$ 或 $q'$ 来指定。于是 $\varPi$ 或 $\varPi'$ 的任何叶都与 $M$ 交于唯一点。构造的 $\SO(3)$ 对称性保证 $M$ 上诱导的复结构（差一符号）就是 $M\cong S^2\cong\bCP^1$ 的自然复结构。但事实上符号相反。看出这点的一个办法是注意簇 $Y$ 有全纯对合 $\zeta:(x,y,z)\to(-x,-y,-z)$。这个对合（它改变全纯辛形式 $\Omega=\d x\d y/z$ 的符号）交换 $\varPi$ 与 $\varPi'$，因为它反转式 (\ref{pok}) 的矩阵 $T$ 的符号，从而交换两个特征值。另一方面，在由实 $x,y,z$ 定义的两球 $M$ 上，映射 $(x,y,z)\to-(x,y,z)$ 反转定向，且反全纯地作用。具体说，$\varPi$ 在 $M$ 上诱导的复结构中 $M$ 由 $q$ 全纯参数化，$\varPi'$ 诱导的复结构中 $M$ 由 $q'$ 全纯参数化。由式 (\ref{pola})，若 $j,x,y,z$ 全为实，则 $q'$ 与 $q$ 满足
\be\label{inverso} q'=-\frac{1}{\bar q}. \ee
故 $\varPi$ 与 $\varPi'$ 诱导的两个复结构相反；一个中的全纯函数在另一个中是反全纯的。

一个有些类似但带有趣转折的例子：仍考虑式 (\ref{nuuxx}) 的 $Y'$，但现在取 $\j^2>0$，使 $s$ 为虚。$\tau$ 的不动点轨迹仍由实 $x,y,z$ 参数化。$\j^2>0$ 的实代数流形 $x^2-y^2-z^2=\j^2$ 有两个分支，$x>0$ 或 $x<0$。分别记作 $M$ 与 $M'$。$M$ 与 $M'$ 各是 $\SO(2,1)/\SO(2)=\SL(2,\R)/\UU(1)$ 的拷贝，也可实现为复上半平面 $\mathcal H$。实 $q$ 或 $q'$ 的 $\varPi$ 或 $\varPi'$ 叶根本不与 $M$ 或 $M'$ 相交。（这由式 (\ref{eipro}) 推出；$\j^2>0$ 时 $s=\pm\sqrt{-\j^2}$ 为虚，故该方程对实 $x,y,z$ 与实 $q$ 或 $q'$ 显然无解。）从 $Y$ 到 $W\cong\bCP^1$ 的投影 $\pi,\pi'$ 把 $M$ 映到 $W$ 中的上半或下半平面，把 $M'$ 映到相反的半平面。$\SO(2,1)$ 对称性保证 $M$ 与 $M'$ 上诱导的复结构（仍差定向选择）等价于 $\mathcal H$ 的标准复结构。

我们讨论的两类全纯极化对量子化有用，后面会看到。很可能某些其他类型也同样有用。例如，几何量子化中可以有介于实极化与复极化之间的极化——粗略地说，在某些方向为实、在另一些方向为复的极化。很可能可以对 $Y$ 的全纯极化加条件，使 $M$ 的膜量子化与用这种中间型极化的几何量子化相联系。
